\documentclass[12pt]{article}
\usepackage[utf8]{inputenc}
\usepackage[T1]{fontenc}
\usepackage{hyperref}
\title{Darang Deng Language Resources for K-12 Curriculum Development}
\author{Compiled by Claude Sonnet 5.5 (claude-sonnet-5-5) for Andr\'as Kornai}
\date{}
\begin{document}
\maketitle

\section{Introduction}
Darang Deng (ISO 639-3 \texttt{dat}) is, per the sources found, an alternate name of the language now coded Digaro Mishmi (\texttt{mhu}), spoken in Arunachal Pradesh, India, and Zayü County, Tibet, China. The code \texttt{dat} was reportedly retired in 2008 as a duplicate. Speaker estimates vary widely (about 6,000 to 45,000).

\section{Primary Educational Resources}
None found in search results.

\section{Linguistic and Reference Materials}
\begin{itemize}
\item Glottolog entry for Darang Deng: \url{https://glottolog.org/resource/languoid/id/dara1255}
\item Wikipedia, Digaro Mishmi language: \url{https://en.wikipedia.org/wiki/Digaro_Mishmi_language}
\item Endangered Languages Project entry (cites van Driem 2007): \url{https://endangeredlanguages.com/node/39962}
\item Omniglot, Digaro: \url{https://mail.omniglot.com/writing/digaro.htm}
\end{itemize}

\section{Religious and Community Resources}
Joshua Project reports Bible translation status as portions only: \url{https://www.joshuaproject.net/index.php/languages/mhu}

\section{Digital and Audio Resources}
None found beyond language overview pages such as Wikitongues: \url{https://wikitongues.org/languages/mhu/}

\section{Curriculum Development Context}
Documentation is thin. The language is usually written in Latin script, sometimes Devanagari. No school use was found; a curriculum would depend on new work with Digaro Mishmi communities.
\end{document}